\section{Three-Point Audit Protocol}
\label{app:verification}

The verification system implements a structured audit for every claim
produced by the system.  This appendix provides formal definitions of
each audit component, specifies their composition into a four-point
verification protocol (three conceptual tests plus Z3 finite-model
checking), and documents the protocol's application to all 34 derivation
steps (Steps 0--33, Phases A--I).

\subsection{Audit Components: Formal Definitions}

\begin{definition}[Hidden Chooser Test (HC)]
\label{def:hc}
Let $s_i$ be derivation step $i$ with input state $I_i$ (the conclusions
of all predecessor steps) and output $O_i$.  The Hidden Chooser Test
verifies that $O_i$ is uniquely determined by $I_i$ and A0:
\[
  \mathrm{HC}(s_i) = \mathrm{PASS} \quad\iff\quad
  \forall\, O' \neq O_i:\;
  O' \text{ is either A0-violating or isomorphic to } O_i.
\]
Formally, ``isomorphic'' means there exists a computable bijection
$\phi \colon O_i \to O'$ such that $\phi$ preserves all quotient-level
properties (truth classes, fiber structure, test outcomes).  A step
\emph{passes HC} if no non-isomorphic A0-compliant alternative exists.
\end{definition}

\begin{definition}[Encoding Invariance Check (EI)]
\label{def:ei}
Let $\Carrier_1$ and $\Carrier_2$ be two carrier representations with
computable bijection $\phi \colon \Carrier_1 \to \Carrier_2$.  Let
$\Ledger_2 = \{(\phi \circ \tau \circ \phi^{-1},\; a) \mid
(\tau, a) \in \Ledger_1\}$ be the induced ledger under $\phi$.  The
Encoding Invariance Check verifies:
\[
  \mathrm{EI}(C, w, \phi) = \mathrm{PASS} \quad\iff\quad
  \TQ{\Ledger_1} \cong \TQ{\Ledger_2},
\]
where $\cong$ denotes isomorphism of quotient structures (preserving
equivalence classes and their separator distances).  This check must hold
for all $\phi \in \gauge$ (the gauge group of Step~12).  EI is the
gauge invariance test applied to the derivation itself: it confirms that
no derivation step depends on a particular bit-level encoding.
\end{definition}

\begin{definition}[Distinguishability Equivalence (DE)]
\label{def:de}
The Distinguishability Equivalence check verifies bidirectional
consistency of A0 application:
\begin{align}
  \text{No false merges:}\quad &
    [x] = [y] \;\Longrightarrow\;
    \nexists\, \tau \in \Del^*:\; \UTM(\tau, x) \neq \UTM(\tau, y).
    \label{eq:no-false-merge} \\
  \text{No false splits:}\quad &
    \bigl(\forall\, \tau \in \Del^*:\; \UTM(\tau, x) = \UTM(\tau, y)\bigr)
    \;\Longrightarrow\; [x] = [y].
    \label{eq:no-false-split}
\end{align}
Equation~\eqref{eq:no-false-merge} ensures that merged objects are
genuinely indistinguishable; equation~\eqref{eq:no-false-split} ensures
that distinguishable objects are never incorrectly identified.  Together
they assert that the truth quotient $\equivL$ is the \emph{exact}
equivalence induced by the test set.
\end{definition}

\begin{definition}[Z3 Finite-Model Verification]
\label{def:z3}
For steps with verifiable algebraic properties (equivalence relations,
gauge group axioms, metric axioms, etc.), Z3 \citep{demoura2008} checks
these properties on finite models:
\begin{itemize}[nosep]
  \item \textbf{Satisfiability checks} (\texttt{sat}): confirm that a
        finite model satisfying the stated properties exists.
  \item \textbf{Unsatisfiability checks} (\texttt{unsat}): confirm that
        no finite model violating the stated properties exists (i.e., the
        property holds for all models in the finite domain).
\end{itemize}
See Appendix~\ref{app:z3} for scripts and results.

\textbf{Honest disclosure:} Satisfiability checks verify existence of a
conforming model, not universal validity.  Unsatisfiability checks are
stronger (universal over the finite domain) but do not extend to
infinite domains without additional mechanization (Lean/Coq).
\end{definition}

\subsection{Composition: Four-Point Audit Protocol}

\begin{definition}[Four-Point Audit]
\label{def:four-point}
A claim $C$ with witness $w$ \emph{passes the four-point audit} if and
only if all four components return PASS:
\begin{align}
  \mathrm{HC}(C, w) &= \mathrm{PASS}
    \tag{Hidden Chooser} \label{eq:audit-hc} \\
  \mathrm{EI}(C, w, \phi) &= \mathrm{PASS}
    \quad \forall\, \phi \in \gauge
    \tag{Encoding Invariance} \label{eq:audit-ei} \\
  \mathrm{DE}(C, w) &= \mathrm{PASS}
    \tag{Distinguishability} \label{eq:audit-de} \\
  \mathrm{Z3}(C) &= \mathrm{SAT} \;\text{or}\; \mathrm{UNSAT}^\dag
    \quad \text{(where applicable)}
    \tag{Finite-Model Check} \label{eq:audit-z3}
\end{align}
where $^\dag$ indicates negation-checking (the negation of the property
is unsatisfiable, proving the property holds universally on the finite
domain).
\end{definition}

\begin{remark}[Protocol Coverage]
The three conceptual tests (HC, EI, DE) are complementary:
\begin{itemize}[nosep]
  \item HC guards against \emph{hidden choices} (non-forced design
        decisions masquerading as theorems).
  \item EI guards against \emph{encoding dependence} (results that change
        under representation switching).
  \item DE guards against \emph{quotient errors} (incorrect merging or
        splitting of descriptions).
\end{itemize}
Together they ensure that each derivation step is forced (HC), invariant
(EI), and faithful (DE).  Z3 provides independent machine-checked
confirmation of algebraic properties on finite models.
\end{remark}

\subsection{Application to Derivation Steps}

\begin{table}[h!]
\centering
\small
\begin{tabular}{@{}clcccl@{}}
\toprule
\textbf{Step} & \textbf{Name} & \textbf{HC} & \textbf{EI} & \textbf{DE} & \textbf{Z3} \\
\midrule
0  & $\noth$ + Output Gate & \checkmark & \checkmark & \checkmark & --- \\
1  & Carrier $\Carrier$, $\FinSlice$ & \checkmark & \checkmark & \checkmark & --- \\
2  & $\SdMap$ syntax & \checkmark & \checkmark & \checkmark & --- \\
3  & Executability & \checkmark & \checkmark & \checkmark & --- \\
4  & $\Del^*$, $\CanEnum$ & \checkmark & \checkmark & \checkmark & --- \\
5  & Ledger $\Ledger$ & $\checkmark^{\ddagger}$ & \checkmark & \checkmark & --- \\
6  & Truth quotient $\TQ{\Ledger}$ & \checkmark & \checkmark & \checkmark & \texttt{sat} \\
7  & $\ObsOpen$, $\Eraser$, sim* & $\checkmark^{\ddagger}$ & \checkmark & \checkmark & --- \\
8  & Entropic time $\Del\Tim$ & \checkmark & \checkmark & \checkmark & \texttt{sat} \\
9  & Entropy, Budget & \checkmark & \checkmark & \checkmark & --- \\
10 & Energy $\EnergyLedger$ & \checkmark & \checkmark & \checkmark & --- \\
11 & $\opnoth$ & \checkmark & \checkmark & \checkmark & --- \\
12 & Gauge $\gauge$ & \checkmark & \checkmark & \checkmark & \texttt{sat} \\
13 & $\Epistemic$/$\Decision$ & \checkmark & \checkmark & \checkmark & --- \\
14 & $\ExactGate$ + Split & \checkmark & \checkmark & \checkmark & \texttt{sat} \\
15 & Myhill--Nerode & \checkmark & \checkmark & \checkmark & \texttt{sat} \\
16 & Separator $\SepDist$ & \checkmark & \checkmark & \checkmark & \texttt{sat} \\
17 & $\Pit$-consistency & \checkmark & \checkmark & \checkmark & --- \\
18 & Bellman minimax & $\checkmark^{\S}$ & \checkmark & \checkmark & \texttt{sat} \\
19 & Self-hosting $\SelfHost$ & $\checkmark^{\P}$ & \checkmark & \checkmark & \texttt{sat} \\
20 & Binary interface & \checkmark & \checkmark & \checkmark & --- \\
21 & Final chain & \checkmark & \checkmark & \checkmark & \texttt{sat} \\
\midrule
22 & Universal closure $\UnivOp$ & \checkmark & \checkmark & \checkmark & --- \\
23 & Closure-defect $\ClosureDefect{Q}$ & \checkmark & \checkmark & \checkmark & --- \\
24 & Consciousness $\ConsProj$ & \checkmark & \checkmark & \checkmark & --- \\
25 & Trit field $\TritField$ & \checkmark & \checkmark & \checkmark & --- \\
\midrule
26 & Continuum limit $\Omega$ & \checkmark & \checkmark & \checkmark & \texttt{sat} \\
27 & Measure $\mu_\Pit$ & \checkmark & \checkmark & \checkmark & \texttt{sat} \\
28 & Dirichlet $\mathcal{E}$ & $\checkmark^{\dagger}$ & \checkmark & \checkmark & \texttt{sat} \\
29 & Semigroup $\{T_t\}$ & \checkmark & \checkmark & \checkmark & \texttt{sat} \\
30 & Metric $g_{ij}$ & \checkmark & \checkmark & \checkmark & \texttt{sat} \\
31 & Complex $J$ & $\checkmark^{\dagger}$ & \checkmark & \checkmark & \texttt{sat} \\
32 & Symplectic $\omega$ & \checkmark & \checkmark & \checkmark & \texttt{sat} \\
33 & K\"{a}hler $(g,J,\omega)$ & \checkmark & \checkmark & \checkmark & \texttt{sat} \\
\bottomrule
\end{tabular}
\caption{Four-point audit results for all 34 derivation steps (Steps 0--33).
$^{\dagger}$Steps~28 and~31 involve modeling commitments (conductance $w=1/d^2$ and smoothness).
$^{\ddagger}$Forced within the paper's pure non-disturbing test model
(commutativity is a modeling commitment; see Appendix~\ref{app:oq:quantum}).
$^{\S}$Forced as the unique non-dominated prior-free policy; Bayesian
alternatives are A0-compatible when a prior is available.
$^{\P}$Forced under the closed-universe assumption (no external verifier);
external verification is A0-compatible but outside the model class.
All steps pass; qualifications are honest disclosures of modeling
commitments, not failures.}
\label{tab:audit}
\end{table}

\subsection{Receipt Chain}

Every system output includes a receipt chain encoded via $\SdMap$:
\begin{enumerate}[nosep]
  \item \textbf{Serialization:} Canonical JSON serialization of the
        output (sorted keys, deterministic formatting).
  \item \textbf{Hashing:} SHA-256 hash of the $\SdMap$-encoded canonical
        JSON, producing a 256-bit fingerprint.
  \item \textbf{Chaining:} Each receipt includes the hash of the previous
        receipt, forming a hash chain: $h_i = \mathrm{SHA256}(h_{i-1}
        \,\|\, \SdMap(\text{output}_i))$.
  \item \textbf{Replay:} Any party can recompute the entire chain from
        the problem specification and verify that the reproduced chain
        matches the original hashes.
\end{enumerate}

The receipt chain ensures \emph{non-repudiation}: the system cannot
silently change its output after the fact.  This is the computational
analogue of the ledger's append-only property (Step~5).  The chain is
\emph{deterministic}: given the same problem specification and budget,
replay produces identical hashes, confirming that the system's behavior
is fully reproducible.

\subsection{Self-Hosting Verification}

The self-hosting property ($\SelfHost = \mathcal{S}$, Step~19) means
the system can verify its own derivation.  The verification procedure
is:
\begin{enumerate}[nosep]
  \item \textbf{Encoding:} Encode each of the 34 derivation steps as a
        compiled finite contract $P_i = \langle q_i, \mathcal{A}_i,
        \mathcal{V}_i, B_i, \pi_i \rangle$, where $q_i$ is the claim,
        $\mathcal{A}_i$ the answer space, $\mathcal{V}_i$ the three-point
        audit, $B_i$ the verification budget, and $\pi_i$ the cost
        assignment.
  \item \textbf{Execution:} Run the system $\mathcal{S}$ on each $P_i$
        using the three-point audit (HC $+$ EI $+$ DE) as the verifier
        $\mathcal{V}_i$.
  \item \textbf{Confirmation:} Verify that all 34 executions (Steps
        0--33) produce $\UNIQUE$ (with modeling-commitment qualifications
        at Steps~5, 18, 19, 28, 31).
  \item \textbf{Certification:} The receipt chain linking all 34
        verifications is the self-hosting certificate.  This certificate
        is itself a finite object in $\Carrier$, verifiable by replay.
\end{enumerate}

\begin{remark}[Scope of Self-Hosting]
Self-hosting is \emph{object-level} verification: the system checks each
derivation step's conclusion given its premises.  It does \emph{not}
verify the meta-claim ``the system is consistent'' or ``the audit
procedure is correct,'' which would violate G\"{o}del's second
incompleteness theorem \citep{godel1931}.  The three-point audit can
verify specific claims but cannot prove the general correctness of the
audit procedure itself.
\end{remark}
